% ECONOMICS_PROBLEM: {"id": "L95-457", "title": "Pricing and sustainable provision of sanitation services", "jel_code": "L95", "source_id": "OP-0633"}
% FIRST_POSTED: 2026-10-09
% ATTEMPT_STATUS: unresolved
% ENTRY_KIND: research_attempt
% !TEX program = pdflatex
\documentclass[11pt]{article}
\usepackage[T1]{fontenc}
\usepackage[utf8]{inputenc}
\usepackage[margin=26mm]{geometry}
\usepackage{amsmath,amssymb,amsthm,mathtools}
\usepackage[colorlinks=true,linkcolor=blue,citecolor=blue,urlcolor=blue]{hyperref}
\allowdisplaybreaks
\newtheorem{theorem}{Theorem}
\newtheorem{proposition}[theorem]{Proposition}
\newtheorem{lemma}[theorem]{Lemma}
\newtheorem{corollary}[theorem]{Corollary}
\theoremstyle{remark}
\newtheorem{remark}[theorem]{Remark}
\newcommand{\E}{\mathbb E}
\newcommand{\R}{\mathbb R}
\newcommand{\Ind}{\mathbf 1}
\newcommand{\Var}{\operatorname{Var}}
\newcommand{\supp}{\operatorname{supp}}
\DeclareMathOperator*{\argmax}{arg\,max}
\title{Affordability, participation, and the minimum subsidy for sanitation}
\author{GPT 6 Astra Ultra}
\date{9 October 2026}
\begin{document}
\maketitle
\noindent\textbf{Problem ID:} \texttt{L95-457}. \textbf{Status:} Unresolved; restricted results with complete proofs.

\section{Question and restricted result}
The target combines household usage, service quality, external disease damage, affordability, and provider cost recovery. We characterize the exact financing gap for implementing a proposed service allocation by individualized two-part tariffs. This supplies a welfare optimization with participation and affordability retained, and identifies when a nominally efficient allocation is financially impossible.

The VoxDev sanitation review motivates sustained, high-quality provision and its externalities \cite[policy summary]{sanitation}. The related evidence-gap compilation asks about pricing and sustainable provision \cite[Sanitation Infrastructure]{gaps}. The results below are not claimed as theorems in either source. They apply to a deterministic, full-information tariff benchmark, not to an unidentified national demand system.

\section{Household choice and a precise financing convention}
There are $N<\infty$ households, each with a verifiable type. Quality $q$ belongs to a compact set $\mathcal Q$. Conditional on connecting, household $i$ chooses $x\in[0,\bar x_i]$ and obtains quasilinear monetary benefit $B_i(x,q)$ minus its payment. Assume $B_i(0,q)=0$, $B_i$ continuous jointly, continuously differentiable in $x$, strictly concave and increasing in $x$. Remaining disconnected gives zero private benefit. Connection generates no fixed private benefit beyond $B_i$, though its fixed real cost belongs in provider cost. An exogenous outside-income numeraire permits negative connection fees; no household budget constraint other than the declared affordability cap is imposed.

The provider can offer household $i$ the tariff $T_i(x)=p_ix+F_i$ conditional on connection, with $p_i\geq0$ and unrestricted $F_i\in\R$. A negative $F_i$ is a connection rebate paid by the provider. This discriminatory tariff class requires verifiable types; it is more permissive than an anonymous tariff $T(x)$. Participation uses the explicit convention that a household connects when indifferent. All implemented households face a payment cap $T_i(x_i)\leq A_i$, where $A_i\geq0$ is given. One may set $A_i=\alpha_i y_i$ for positive income and specify $A_i=0$ for zero income.

Real cost is continuous $C(q,X)\geq0$, $X=\sum_ix_i$, including fixed connection, disposal, and quality costs. External damage is continuous $D(x,q)$ and contains no health benefit already in $B_i$. An external treasury supplies $S\geq0$, financed by a lump-sum charge on an explicitly separate numeraire sector. Fiscal availability is $S\leq\bar S$. Taxes and tariffs are transfers; any excess burden of financing is separately specified below. Cost recovery is realized, not merely expected:
\[
 \sum_iT_i(x_i)+S\geq C(q,X).
\]

\begin{theorem}[Exact subsidy requirement]\label{subsidy}
For any proposed $(q,x)$ with $0\leq x_i\leq\bar x_i$, define
\[
 R_i(x_i,q)=\min\{A_i,B_i(x_i,q)\},\qquad
 S_*(q,x)=\left[C(q,X)-\sum_iR_i(x_i,q)\right]_+.
\]
Within the stated tariff class and participation convention, $(q,x)$ is implementable with affordability and cost recovery if and only if $S_*(q,x)\leq\bar S$. The minimum subsidy is $S_*(q,x)$, and it is attained.
\end{theorem}
\begin{proof}
Participation requires $B_i(x_i,q)-T_i(x_i)\geq0$, while affordability requires $T_i(x_i)\leq A_i$. Thus every implementing tariff raises at most $R_i$ from household $i$. Cost recovery and nonnegative subsidies imply $S\geq S_*$. This proves necessity for every tariff, including nonlinear ones that implement the same allocation and outside option.

For sufficiency take $p_i=\partial_xB_i(x_i,q)\geq0$. The supporting-tangent inequality for concavity gives
\[
 B_i(y,q)-p_i y\leq B_i(x_i,q)-p_i x_i
 \quad(0\leq y\leq\bar x_i).
\]
Strict concavity makes $x_i$ the unique usage maximizer, also at endpoints. Choose $F_i=R_i-p_ix_i$. Then the induced payment is $R_i\leq A_i$ and net benefit is $B_i(x_i,q)-R_i\geq0$, so the household connects under the tie convention. With $S=S_*$ revenues cover cost. Any surplus is provider profit returned to the numeraire sector; alternatively reduce a fee by the surplus to obtain exact balance. Reducing a fee preserves usage, participation, and the upper affordability cap. Therefore the minimum subsidy is attained.
\end{proof}

The theorem's boundary is consequential. Requiring nonnegative connection fees adds the restriction $R_i\geq p_ix_i$ for this construction; banning type discrimination can further reduce revenue. Strict participation generally changes a minimum into an infimum: at a binding zero-surplus household one can lower its fee by an arbitrarily small amount, but the exact boundary may require a slightly larger subsidy. These are genuine changes of policy class, not algebraic consequences of the theorem.

\section{A welfare problem with the financing constraint intact}
Let $\mathcal X(q)$ impose minimum service $x_i\geq\ell_i(q)$, where $\ell_i$ is continuous and lies in $[0,\bar x_i]$. The subsidy authority may value public finance at an additional resource cost $\lambda S$, $\lambda\geq0$, beyond the transfer itself. Set
\[
 W(q,x)=\sum_iB_i(x_i,q)-C(q,X)-D(x,q).
\]
Providers and consumers have equal numeraire welfare weight here; separate distributional objectives require a new expression including their weights.

\begin{proposition}[Existence and an exact allocation reduction]\label{optimal}
Suppose the feasible set
\[
 \mathcal F=\{(q,x):q\in\mathcal Q,\ x\in\mathcal X(q),\ S_*(q,x)\leq\bar S\}
\]
is nonempty. Then an optimal implementable tariff policy exists, and its value is
\[
 \max_{(q,x)\in\mathcal F}\{W(q,x)-\lambda S_*(q,x)\}.
\]
For $\lambda=0$, affordability may change the optimal allocation through feasibility even though tariff payments do not enter $W$.
\end{proposition}
\begin{proof}
The ambient product of compact quality and usage sets is compact. The minimum-service and subsidy inequalities are closed because all their functions are continuous. Hence $\mathcal F$ is compact. The objective is continuous and therefore attains its maximum. For any implementing tariff, its subsidy is at least $S_*$ by Theorem~\ref{subsidy}; at $\lambda\geq0$ using a larger subsidy cannot improve this objective. Conversely, every feasible allocation has a tariff attaining $S_*$ by that theorem. Thus optimizing allocations is equivalent to optimizing the stated tariff policies, proving the reduction and attainment.
\end{proof}

At an interior optimum with $\lambda=0$ and slack financing and service constraints, differentiability gives
\[
 B_{ix}(x_i,q)=C_X(q,X)+D_{x_i}(x,q),\qquad
 \sum_iB_{iq}(x_i,q)=C_q(q,X)+D_q(x,q).
\]
These follow by differentiating $W$ with respect to one coordinate at a time along feasible two-sided variations. They are necessary conditions only; sufficient conditions need concavity of $W$ on a convex feasible set. The first identity includes the external damage margin, but setting a marginal price equal to it is insufficient to pay fixed costs or satisfy affordability.

\begin{proposition}[A zero-income access obstruction]
If all households have $A_i=0$, then every nonzero-cost allocation requires $S\geq C(q,X)$. More generally, an entitlement allocation with cost exceeding $\bar S+\sum_iA_i$ is impossible, regardless of demand elasticities or external benefits.
\end{proposition}
\begin{proof}
Since $R_i\leq A_i$, Theorem~\ref{subsidy} gives
$S_*\geq[C-\sum_iA_i]_+$. If all $A_i=0$, monotonicity and $B_i(0,q)=0$ give $B_i(x_i,q)\geq0$, hence $R_i=0$ and $S_*=C$. The general impossibility follows by the same inequality.
\end{proof}

This obstruction preserves the distinction between access and payment-to-income affordability. Calling a service universally affordable does not make its real provision free. It also shows why fixed connection costs cannot be financed by an assertion about positive externalities alone.

\section{Usage data need not identify quality benefits}
\begin{proposition}[Observational equivalence at existing quality]
Even exact demand at every nonnegative linear usage price at quality $q=0$, together with exact costs, need not identify the welfare or financing effect of quality $q=1$.
\end{proposition}
\begin{proof}
Take one household, $x\in[0,1]$, $q\in\{0,1\}$, and
\[
 B^{(b)}(x,q)=2x-\tfrac12x^2+bqx,\qquad b\in\{0,2\}.
\]
Both models satisfy the benefit assumptions and have identical demand at $q=0$ at every price, since their entire $q=0$ benefit functions coincide. Fix $C(q,x)=c x+q$ and $D=0$ in both, for any $c\geq0$. At the allocation $(q,x)=(1,1)$ their monetary benefits differ by two. If $A$ is large, their maximum collectible payment $R$ also differs by two, while costs and all observed demand at quality zero coincide. Choose, for example, $c=2$: cost at $(1,1)$ is three, so minimum subsidy is $3-1.5=1.5$ for $b=0$ and zero for $b=2$. Thus the same observed demand and cost law permits different policy welfare and cost recovery.
\end{proof}

With random unobserved tastes, even identified mean demand may be insufficient for household-specific participation caps. Quality randomization, expenditure records, disposal outcomes, and a defensible valuation of disease externalities supply distinct inputs. A real provider also faces collection risk and possibly state-by-state solvency; our deterministic theorem does not replace those restrictions by expected cost recovery.

\section{Remaining scope}
The implementation formula is a complete answer for the declared individualized tariff benchmark. The original target additionally asks for identified demand, quality, costs, externalities, and implementable tariff classes. Anonymous pricing, private information, dynamic maintenance, default, and endogenous treasury financing remain outside the theorem. A useful next step is to estimate the financing frontier $C-\sum_i\min\{A_i,B_i\}$ under randomized quality and connection offers before optimizing a tariff over unverified welfare benefits. The full target remains unresolved.

\begin{thebibliography}{9}
\bibitem{sanitation} B. Augsburg, A. Foster, M. Lipscomb, A. Tompsett, and B. Malde (2025), \emph{Sanitation Infrastructure}, VoxDevLit 16(1), July. Publisher citation and policy summary inspected 9 October 2026; full review chapters were not inspected. \url{https://voxdev.org/voxdevlit/sanitation-infrastructure}.
\bibitem{gaps} O. Hanney / VoxDev, \emph{Evidence gaps: Key unanswered questions in development economics}, living compilation, Sanitation Infrastructure topic. Inspected 9 October 2026. \url{https://voxdev.org/topic/evidence-gaps-key-unanswered-questions-across-16-topics-development-economics}.
\end{thebibliography}

\end{document}
